% ============================================================
%  第 1 章 Introduction 导读（体例：中文导读 + 关键句精读 + 原文定位）
% ============================================================
\chapter{Introduction}
\begin{center}
\begin{minipage}{0.9\textwidth}
\itshape\small\color{dsgray}
本章确立全书的基本概念：分布式系统与去中心化系统的区别、为什么“分散”
本身不是目标、六大设计目标（资源共享、分布透明、开放性、可靠性、安全、
可扩展性），以及三类系统（高性能计算、信息处理、普适计算）。
最后给出著名的“八个错误假设”，它们是全书反复要打破的思维惯性。
\end{minipage}
\end{center}

\sechead{1.1 From networked systems to distributed systems}{从网络系统到分布式系统}

计算机从 1945 年到 1985 年之间大多是彼此孤立的大型机。转折出现在 1980 年代中
期，两股力量改变了局面：\textbf{强大的微处理器}（从 8 位到 16/32/64 位再到多核）
与\textbf{高速计算机网络}（LAN 与 WAN）。叠加设备小型化（智能手机、树莓派这类
“纳米计算机”），联网计算迅速渗入汽车、电网、桥梁等基础设施——这也让它们变得
“可被攻击”。

\subsection*{1.1.1 Distributed versus decentralized systems}

作者给出的区分，关键在“必要（necessity）”与“充分（sufficiency）”之分：

\begin{itemize}
  \item \textbf{去中心化（decentralized）}：\emph{被迫}分散。原因往往是行政边界、
        信任缺失或地理约束。三个典型例子：联邦学习（数据不能出组织边界，
        于是把训练送到数据那里）、区块链/分布式账本（参与方互不信任，
        交易必须公开可验证）、地理上天然分散的监控系统（卫星、交通）。
  \item \textbf{分布式（distributed）}：\emph{主动}分散，以换取可扩展性与容错。
        典型例子：Gmail（逻辑上只有两台服务器，背后是海量机器）、CDN（Akamai，
        2022 年超 40 万台服务器）、家用 NAS（多用户共享文件的迷你分布式系统）。
\end{itemize}

\begin{closeread}
\pair{A \emph{decentralized system} is a networked computer system in which processes and resources are \emph{necessarily} spread across multiple computers.}{一个\textbf{去中心化系统}，是进程与资源\emph{必然}分散在多个计算机上的网络系统。}
\pair{A \emph{distributed system} is a networked computer system in which processes and resources are \emph{sufficiently} spread across multiple computers.}{一个\textbf{分布式系统}，是进程与资源\emph{充分}分散在多个计算机上的网络系统。}
\end{closeread}

\origin{§1.1.1，原书 p.4}

\begin{ideabox}[Note 1.1：中心化就一定差吗？]
常见误解：中心化无法扩展、且必然带来单点故障。作者强调必须区分
\textbf{逻辑设计}与\textbf{物理设计}。DNS 逻辑上是一棵树，物理上根节点由
13 台根服务器实现、每台又是集群；它高度副本化，几乎不可能被单点打垮。
中心化反而因为“只有一个故障点”而更容易管理、加固，云（cloud）就是高度可扩展
又稳健的中心化方案。
\end{ideabox}

\subsection*{1.1.2 Why the distinction is relevant}

中心化方案通常更简单。把服务拆散是一个需要\textbf{慎重权衡}的决定，因为
分布式/去中心化系统天生困难：存在大量\textbf{非预期依赖}；几乎持续处于
\textbf{部分失败（partial failure）}状态；节点不断来去、系统高度动态；
跨网络跨域使其\textbf{易受攻击}。

\commentary{\textbf{全书立场}：“分散”永远不能是目的本身，原则是\emph{能少分散就少分散}
（追求“充分”而非“必然”）；能简单就简单，因此本书基本不讨论优化——作者认为
优化带来的复杂度代价高于其性能收益。}

\begin{closeread}
\pair{Throughout this book, we take the standpoint that decentralization can never be a goal in itself.}{贯穿本书的立场是：去中心化永远不能成为目的本身。}
\pair{In principle, the less spreading, the better.}{原则上，分散得越少越好。}
\end{closeread}

\origin{§1.1.2，原书 p.8}

\subsection*{1.1.3 Studying distributed systems}

由于依赖关系无处不在，本书选择从一组\textbf{有限的、互不相同的视角}切入，
每章一个视角：架构、进程、通信、协调、命名、一致性与副本、容错、安全。

\begin{closeread}
\pair{For example, there is no such thing as a separate communication module, or a separate security module.}{例如，并不存在一个独立的“通信模块”，也不存在一个独立的“安全模块”。}
\pair{Each perspective is considered in a separate chapter.}{每一个视角都在单独的一章中讨论。}
\end{closeread}

\origin{§1.1.3，原书 p.8–10}

\sechead{1.2 Design goals}{设计目标}

\subsection*{1.2.1 资源共享（Resource sharing）}
目标：让用户与应用\textbf{方便地访问和共享远程资源}（外设、存储、数据、文件、
服务、网络）。动机包括经济性（共享一套高端存储比人人各买一套便宜）、
协作（协同编辑/视频会议、跨国软件外包）、P2P 文件共享（BitTorrent）。
共享已高度“无感”：第三方维护的共享文件夹，用起来几乎和本机文件夹无异。

\origin{§1.2.1，原书 p.10–11}

\subsection*{1.2.2 分布透明（Distribution transparency）}
目标：\textbf{隐藏}进程与资源物理上分散在多地这一事实。实现方式是
\textbf{中间件（middleware）}，它为应用提供处处一致的接口，把“在哪、怎么访问”
藏起来。原书（引 ISO 1995）列出七种透明性：\textbf{访问、位置、重定位、迁移、
复制、并发、失败}。其中失败透明的核心困难是：\textbf{无法区分“死进程”与
“只是很慢的进程”}。

\begin{closeread}
\pair{An important goal of a distributed system is to hide the fact that its processes and resources are physically distributed across multiple computers.}{分布式系统的一个重要目标，是隐藏“其进程与资源在物理上分散于多台计算机”这一事实。}
\pair{The price for achieving full transparency may be surprisingly high.}{实现完全透明所付出的代价，可能高得惊人。}
\end{closeread}

\origin{§1.2.2，原书 p.11–15（透明性表见图 1.3）}

\begin{ideabox}[关键权衡：透明到什么程度？]
透明并非越多越好。跨时区的“早晨报纸”不该假装还按你本地作息；旧金山到
阿姆斯特丹的往返延迟有物理下限（光速约 35ms 起，实测几百毫秒）；“跨大洲永远一致”
的复制会让单次更新耗时数秒，藏不住。作者主张：透明应和性能、可理解性一起权衡，
\textbf{有时显式暴露分布反而更好}（如基于位置的服务）。Waldo 等（1997）指出，
用 RPC 掩盖分布会导致难以理解的语义，因为过程调用一旦跨失败链路，性质就变了。
\end{ideabox}

\subsection*{1.2.3 开放性（Openness）}
开放的分布式系统：其组件能\textbf{被别的系统方便地使用/集成}，自身也常由“来自
别处”的组件构成。接口通常用 \textbf{IDL} 描述，能精确刻画\textbf{语法}，但
\textbf{语义}往往只能用自然语言非正式描述。理想规范应\textbf{完整}且\textbf{中立}，
从而带来\textbf{互操作性（interoperability）}、\textbf{可移植性（portability）}、
\textbf{可扩展性（extensibility）}。实现灵活性的关键是\textbf{策略（policy）
与机制（mechanism）分离}。

\begin{closeread}
\pair{To be open means that components should adhere to standard rules that describe the syntax and semantics of what those components have to offer.}{所谓开放，意味着组件应当遵循统一的规则，这些规则描述了组件所提供的语法与语义。}
\pair{What we need is a separation between policy and mechanism.}{我们需要的，是把策略与机制分离开来。}
\end{closeread}

\origin{§1.2.3，原书 p.15–18（Note 1.4 讨论分离过度的代价）}

\subsection*{1.2.4 可靠性（Dependability）}
可靠性指系统\textbf{被信赖按预期运行的程度}。分布式系统中由于部分失败，
应假设\textbf{任意时刻总有部分失败在发生}，核心目标是掩盖失败及其恢复，即
\textbf{容错（fault tolerance）}。要求包括\textbf{可用性、可靠性、安全性
（safety）、可维护性}。故障三层次：\emph{fault}（成因）→ \emph{error}
（可能引发失败的状态）→ \emph{failure}（无法履约）。故障按时序分
\textbf{瞬时/间歇/永久}。

\begin{closeread}
\pair{Availability is defined as the property that a system is ready to be used immediately.}{可用性被定义为：系统随时处于可供立即使用的状态。}
\pair{Reliability refers to the property that a system can run continuously without failure.}{可靠性指的是：系统能够持续无故障地运行。}
\end{closeread}

\origin{§1.2.4，原书 p.18–21}

\begin{ideabox}[可用性 $\neq$ 可靠性]
每小时随机宕机 1 毫秒的系统，可用性超过 99.9999\%，却\textbf{不可靠}；
从不崩溃、但每年 8 月停机两周的系统，可靠性高，可用性只有 96\%。
常用度量：MTTF（平均无故障时间）、MTTR（平均修复时间）、MTBF = MTTF + MTTR。
\end{ideabox}

\subsection*{1.2.5 安全（Security）}
“不安全的分布式系统就不可靠。”核心是\textbf{认证（authentication，验证所声称
的身份）}与\textbf{信任（trust）}。本节只引入最基本密码学工具（详见第 9 章）：
\textbf{对称密码}（加解密同一密钥）、\textbf{非对称/公钥密码}（公钥加密、
私钥解密；反向使用即\textbf{数字签名}）、\textbf{哈希函数}（定长、改一点就变、
不可反推）。

\begin{closeread}
\pair{In a symmetric cryptosystem, encryption and decryption takes place with a single key.}{在对称密码系统中，加密与解密使用同一把密钥。}
\pair{In an asymmetric cryptosystem, the keys used for encryption and decryption are different.}{在非对称密码系统中，用于加密和解密的密钥是不同的。}
\end{closeread}

\origin{§1.2.5，原书 p.21–23（记号表见图 1.4）}

\subsection*{1.2.6 可扩展性（Scalability）}
按三个维度衡量（Neuman 1994）：\textbf{规模（size）}、\textbf{地理（geographical）}、
\textbf{行政（administrative）}。规模瓶颈的根因是 CPU、存储与 I/O、网络；
地理的难点是同步通信延迟、可靠性、带宽与缺乏多点通信；行政最棘手，涉及政策冲突。
向外扩展（scale out）只有三种手段：\textbf{隐藏通信延迟}、\textbf{分区与分布}、
\textbf{复制}。复制（含\textbf{缓存}，区别在于缓存由\emph{使用者}而非\emph{所有者}
决定）的代价是\textbf{一致性}。

\begin{closeread}
\pair{A system can be scalable regarding its size, meaning that we can easily add more users and resources to the system without any noticeable loss of performance.}{系统在规模上可扩展，意味着我们可以轻松地向系统加入更多用户与资源，而性能没有可察觉的下降。}
\pair{There are basically only three techniques we can apply: hiding communication latencies, distribution of work, and replication.}{我们基本上只能采用三种技术：隐藏通信延迟、分摊工作、以及复制。}
\end{closeread}

\origin{§1.2.6，原书 p.24–32（Note 1.5 用排队论分析规模可扩展性）}

\sechead{1.3 A simple classification}{分布式系统的简单分类}

\subsection*{1.3.1 高性能分布式计算}
\textbf{集群计算（cluster）}：硬件\textbf{同构}，一批相似节点经高速互连，用于
并行编程；现流行\textbf{多内核（multikernel）}——满配 OS 与轻量内核并存。
\textbf{网格计算（grid）}：\textbf{异构}，把不同组织的资源聚成\textbf{虚拟组织
（virtual organization）}。Foster 等的四层架构：fabric / connectivity / resource /
collective，上接 application；中间三层即“网格中间件”。

\origin{§1.3.1，原书 p.32–37（网格四层见图 1.10）}

\subsection*{1.3.2 分布式信息系统}
\textbf{分布式事务处理}：事务满足 \textbf{ACID}（原子、一致、隔离、持久）。
分布式事务常拆成\textbf{嵌套事务}：即使某子事务已提交、其父事务后来中止，
该子事务结果也必须回滚——所以“持久性”只对\textbf{顶层事务}成立。核心组件是
\textbf{TP monitor}，按\textbf{分布式提交}协议协调子事务提交。
\textbf{企业应用集成（EAI）}的常见方式：\textbf{RPC}、\textbf{RMI}、
\textbf{MOM（发布-订阅）}。

\begin{closeread}
\pair{The characteristic feature of a transaction is either all of these operations are executed or none are executed.}{事务的特征是：这些操作要么全部执行，要么全部不执行。}
\pair{A transaction-processing monitor's main task is to allow an application to access multiple server/databases by offering it a transactional programming model.}{事务处理监视器（TP monitor）的主要任务，是通过提供一个事务式编程模型，让应用得以访问多个服务器/数据库。}
\end{closeread}

\origin{§1.3.2，原书 p.37–43（Note 1.7 讨论四种集成方式）}

\subsection*{1.3.3 普适系统（Pervasive systems）}
因移动/嵌入式设备而不同于前两类：\textbf{用户与组件的边界模糊}，没有单一
屏/键接口，而是大量\textbf{传感器}与\textbf{执行器}。三类（有重叠）：

\begin{itemize}
  \item \textbf{泛在计算（ubiquitous）}：系统\textbf{持续存在}，交互多为“隐式输入”。
        Poslad（2009）五项要求：分布、交互、情境感知、自治、智能。
  \item \textbf{移动计算（mobile）}：设备多样且\textbf{位置随时间变化}；MANET 未流行，
        移动计算回归“移动设备连到服务器”，由此引出 \textbf{MEC（移动边缘计算）}。
  \item \textbf{传感器网络（sensor networks）}：数十至上千小节点，资源与能耗受限；
        关键是把网络当作\textbf{一个数据库}（TinyDB 声明式接口），做\textbf{网内处理
        （in-network processing）}——沿树上传并在汇合处聚合。
\end{itemize}

\noindent 由 IoT 与云连接，形成 \textbf{云–边–物}层级（图 1.17）：越高层越可靠、
容量越大、性能越强；越低层越易实现位置相关能力与低延迟，设备数量也越多。

\origin{§1.3.3，原书 p.43–52（云-边-物层级见图 1.17）}

\sechead{1.4 Pitfalls}{陷阱：八个错误假设}

Peter Deutsch（当年在 Sun）总结了人们初次开发分布式应用时的\textbf{八个错误假设}：

\begin{enumerate}[leftmargin=1.6em]
  \item The network is reliable（网络可靠）
  \item The network is secure（网络安全）
  \item The network is homogeneous（网络同构）
  \item The topology does not change（拓扑不变）
  \item Latency is zero（延迟为零）
  \item Bandwidth is infinite（带宽无限）
  \item Transport cost is zero（传输零成本）
  \item There is one administrator（只有一个管理员）
\end{enumerate}

\begin{ideabox}[为什么这八条是全书索引]
它们正好对应分布式系统独有的性质：可靠性、安全性、异构性、拓扑；延迟与带宽；
传输成本；行政域。本书几乎每个原理，都是在反驳其中某条假设所产生的问题——
例如“可靠网络不存在”导致失败透明不可能；“网络本就不安全”单独成章；
复制是为解决延迟与带宽问题；管理问题穿插在各处。
\end{ideabox}

\origin{§1.4，原书 p.52–53}

\sechead{1.5 Summary}{小结}

\begin{itemize}
  \item 分布式系统 = 进程与资源\textbf{充分}分散；去中心化系统 = \textbf{必然}分散。
        分散本身不是目标，多数分布式选择是为提升可靠性、可扩展性、效率；
        但中心化系统往往更好管理，做决定前要三思——有时又别无选择。
  \item 设计目标：共享资源、保证开放、日益重要的安全，以及隐藏分布细节。
        透明有性能代价且\textbf{永远无法完全实现}；它与可扩展性常相互冲突。
  \item 系统分类：面向\textbf{计算}（集群→网格→云）、面向\textbf{信息处理}、
        面向\textbf{普适}（移动、传感器、泛在）。
  \item 反复出现的陷阱：开发者对底层网络作\textbf{根本错误}的假设，
        事后很难掩盖由此产生的不良行为。
\end{itemize}

\origin{§1.5，原书 p.53–54}

\begin{actions}
  \item 读原书第 1 章对应小节，重点看图 1.1（三种组织）、图 1.3（透明性表）、图 1.10（网格四层）、图 1.17（云-边-物层级）。
  \item 用“必然 vs 充分”重新审视一个你熟悉的系统（如微信/网盘）：哪些分散是被迫的？哪些是主动的？
  \item 把你正在做或想做的项目，逐条对照八个错误假设，写下你默认了哪几条。
  \item 记住三个“反直觉结论”：中心化不必差、透明不是越多越好、可用性 $\neq$ 可靠性。
\end{actions}
